\SourceLeaf{106}
\begin{center}\textbf{[DREI FRAGMENTE ÜBER ELLIPTISCHE MODULFUNCTIONEN.]}\end{center}
\begin{center}\rule{1.5cm}{.4pt}\end{center}
\begin{center}[1.]\end{center}

Ist
\[
\frac{n}{m}=\frac{\mu(1-2e^{-M\pi}+2e^{-4M\pi}-2e^{-9M\pi}+\cdots)^2}{\mu(1+2e^{-M\pi}+2e^{-4M\pi}+2e^{-9M\pi}+\cdots)^2},
\]
so kann man statt $M$ setzen
\[
\frac{1}{2ai}\StepDown{.75}{+\frac{1}{2bi}}\StepDown{1.5}{+\frac{1}{2ci}}\StepDown{2.25}{+\frac{1}{2di}}\StepDown{3}{+\frac{1}{2ei}}\StepDown{3.75}{+\text{etc.}+\frac{1}{M},}
\]
wo $a,b,c,d,e$ u.\,s.\,w. eine beliebige ungerade Menge ganzer reeller Zahlen bedeuten; oder auch
\[
\frac{pM+2qi}{r+2sMi}=\frac{M+2(qr-psMM)i}{rr+4ssMM},
\]
wo $p,q,r,s$ beliebige, der Bedingung
\[
pr+4qs=1
\]
Genüge leistende ganze reelle Zahlen sind.

\begin{center}\rule{3cm}{.4pt}\end{center}
\SourceLeaf{107}
\begin{center}[2.]\end{center}
\begin{center}\textbf{DIE REDUCTION VON $pM$, $qM$, $rM$ AUF DIE EINFACHSTE FORM.}\end{center}

Es sei $M=\dfrac{\alpha+\beta i}{\delta-\gamma i}$, wo $\alpha$, $\beta$, $\gamma$, $\delta$ ganze reelle Zahlen und Zähler und Nenner ohne gemeinschaftlichen Factor. Man setze
\[
\begin{gathered}
\alpha\alpha+\beta\beta=A,\qquad \alpha\gamma+\beta\delta=B,\qquad \gamma\gamma+\delta\delta=C,\\
\alpha\delta-\beta\gamma=\BareRoot{(AC-BB)}=D.
\end{gathered}
\]
Man suche die einfachste Form des Determinanten $-DD$, welche der Form $(A,B,C)$ aequivalent ist; sie sei $(a,b,c)$.

Dann lassen sich die Functionen von $M$ auf Functionen von
\[
\frac{D+bi}{a}
\]
zurückführen. Der Algorithmus ist dieser
\[
\begin{array}{ccc}
\dfrac{D+Bi}{A}=M,&DD+BB=AA',&B+B'=hA',\\[4pt]
\dfrac{D+B'i}{A'}=M',&DD+B'B'=A'A'',&B'+B''=h'A'',\\[4pt]
\dfrac{D+B''i}{A''}=M'',&DD+B''B''=A''A''',&B''+B'''=h''A''',
\end{array}
\]
\begin{center}.\quad .\quad .\quad .\quad .\quad .\quad .\quad .\quad .\quad .\quad .\quad .\quad .\quad .\quad .\quad .\quad .\quad .\quad .\quad .\quad .\end{center}
\[
\begin{array}{rl}
\BareRoot{\dfrac{D+Bi}{A}}\cdot pM=pM'&\text{für gerades }h,\\[4pt]
{}=qM'&\text{für ungerades }h,\\[5pt]\hline\\[-6pt]
\varepsilon^h\BareRoot{\dfrac{D+Bi}{A}}\cdot qM=rM',&\varepsilon=\sqrt{i},\\[5pt]\hline\\[-6pt]
\BareRoot{\dfrac{D+Bi}{A}}\cdot rM=qM'&\text{für gerades }h,\\[4pt]
{}=pM'&\text{für ungerades }h.
\end{array}
\]

Wenn man aus $h$, $h'$, $h''$ u.\,s.\,w. die Transformation von $(A,B,A')$ in $(a,b,c)$ ableitet, so werden deren Elemente (ob sie gerade oder ungerade sind) entscheiden, welche Function von $\dfrac{D+bi}{a}$ mit der gegebnen von $M$ so zusammen\SourceLeaf{108}hängt, dass letztere in
\[
\varepsilon^H\BareRoot{\left(\dfrac{D+Bi}{A}\cdot\dfrac{D+B'i}{A'}\cdot\dfrac{D+B''i}{A''}\cdots\right)}
\]
multiplicirt werden muss.

Wo $M$ nicht rational ist, mag man $D=-1$ setzen und den Algorithmus ebenso bilden; nemlich, wenn $M=g+hi$, so geht man von der Form $\left(\dfrac1g,\dfrac hg,\dfrac{gg+hh}{g}\right)$ (Det. $-1$) aus, sucht ihre Aequivalente etc.

\begin{center}\rule{3cm}{.4pt}\end{center}
\begin{center}[3.]\end{center}
\[
\begin{aligned}
pt&=1+2e^{-\pi t}+2e^{-4\pi t}+2e^{-9\pi t}+\text{etc.}\\
qt&=1-2e^{-\pi t}+2e^{-4\pi t}-2e^{-9\pi t}+\text{etc.}\\
rt&=2e^{-\frac14\pi t}+2e^{-\frac94\pi t}+2e^{-\frac{25}4\pi t}+\text{etc.}
\end{aligned}
\]

Um die Gleichung
\[
\frac{qt}{pt}=A
\]
aufzulösen, setze man $AA=\dfrac nm$ und suche das a.-g. Mittel zwischen $m$ und $n$; es sei dasselbe $=\mu$. Man suche ferner das a.-g. Mittel zwischen $m$ und $\BareRoot{(mm-nn)}$ oder, was dasselbe ist, zwischen $(m+n)$ und $(m-n)$; dieses sei $=\lambda$. Man hat dann $t=\dfrac\mu\lambda$.

Man erhält so nur Einen Werth von $t$; sämmtliche andere werden dann in der Formel
\[
\frac{\alpha t-2\beta i}{\delta-2\gamma ti}
\]
enthalten sein, wo $\alpha$, $\beta$, $\gamma$, $\delta$ alle ganzen Zahlen bedeuten, die der Gleichung:
\[
\alpha\delta-4\beta\gamma=1
\]
Genüge leisten.

Um aus $A$ abzuleiten
\[
B=\dfrac{q\!\left(\frac13t\right)}{p\!\left(\frac13t\right)},
\]
\SourceLeaf{109}ist eine biquadratische Gleichung aufzulösen:
\[
(B-A)^4=4(A-A^3)(B-B^3),
\]
oder
\[
(1-AB)^4=(1-A^4)(1-B^4).
\]

Den vier Wurzeln correspondiren $\frac13t$, $\frac13t+\frac23i$, $\frac13t+\frac43i$, $3t$.

Für $A=\frac12$ suche man zwei a.-g. Mittel
\[
\begin{aligned}
m&=4, & n&=1\\
\mu&=2{,}2430340, &\lambda&=3{,}9364917.
\end{aligned}
\]
Also
\[
\begin{array}{rl}
t=&0{,}56983,\\
\log t=&9{,}7557537,\\
\log(\pi\log e)=&0{,}1349342,\\\hline
&9{,}8906879,\\
\log e^{-\pi t}=&-0{,}7774777,\\
{}=&9{,}2225223,\\
pt=&1{,}33540375,\\
qt=&0{,}66770187.
\end{array}
\]
\begin{center}\rule{3cm}{.4pt}\end{center}
